%%%
%%% Nicola Pero
%%%
%%%
%%%
%%%
%%%

\documentclass[a4paper]{article}
\begin{document}
\author{Nicola Pero nicola.pero@meta-innovation.com}
\title{First Steps in GNUstep GUI Programming (2): NSWindow, NSButton}
\date{First Edition: July 2000; Last Updated: January 2008}
\maketitle

In this tutorial you will learn how to create a window in your
application.  We will put then a button in the window.

\section{Points, Sizes and Rectangles}
The GNUstep base library defines some useful structs for dealing with
two dimensional geometry: \texttt{NSPoint}, \texttt{NSSize} and
\texttt{NSRect}.  It is worth to quickly review them here, before
beginning.

\subsection{NSPoint}
A \texttt{NSPoint} is defined as a struct with two members, the
\emph{x} and \emph{y} coordinate of the point:
\begin{verbatim}
typedef struct _NSPoint
{
  float x;
  float y;
} NSPoint;
\end{verbatim}
So, to access the \texttt{x} and \texttt{y} coordinates of a
\texttt{NSPoint} called \texttt{myPoint}, you just do as in
\texttt{myPoint.x} and \texttt{myPoint.y}.  Please note that the
coordinates of a \texttt{NSPoint} are floats; so, they might be
negative and/or fractionary.

To create a point with given \texttt{x} and \texttt{y} coordinates,
you can use the function (macro) \texttt{NSMakePoint ()}, as in the
following example:
\begin{verbatim}
NSPoint testPoint;

testPoint = NSMakePoint (10, 20);
NSLog (@"x coordinate: %f", testPoint.x); // 10
NSLog (@"y coordinate: %f", testPoint.y); // 20
\end{verbatim}

It might be worth quoting the (predefined) constant
\texttt{NSZeroPoint}, which is a point with zero x coordinate and zero
y coordinate.

\subsection{NSSize}
An \texttt{NSSize} is a struct describing the size of a rectangle,
regardless of its position.
\begin{verbatim}
typedef struct _NSSize
{
  float width;
  float height;
} NSSize;
\end{verbatim}
Using \texttt{NSSize} is completely analogous to using \texttt{NSPoint}, 
as in the following code example: 
\begin{verbatim}
NSSize testSize;

testSize = NSMakeSize (0.5, 51);
NSLog (@"x coordinate: %f", testSize.width); // 0.5
NSLog (@"y coordinate: %f", testSize.height); // 51
\end{verbatim}
\texttt{NSZeroSize} is a constant equal to a size with zero width 
and zero height.

\subsection{NSRect}
An \texttt{NSRect} is a struct describing both the position and the
size of a rectangle:
\begin{verbatim}
typedef struct _NSRect
{
  NSPoint origin;
  NSSize size;
} NSRect;
\end{verbatim}
Using \texttt{NSRect} is similar to using \texttt{NSPoint} and \texttt{NSSize}: 
\begin{verbatim}
NSRect testRect;

testRect = NSMakeRect (8.1, -3, 10, 15);
NSLog (@"x origin: %f", testRect.origin.x); // 8.1
NSLog (@"y origin: %f", testRect.origin.y); // -3
NSLog (@"width: %f", testRect.size.width); // 10
NSLog (@"height: %f", testRect.size.height); // 15
\end{verbatim}
Note that you first access the \texttt{origin}, and then its coordinates, 
and similarly for the \texttt{size}. 

The constant \texttt{NSZeroRect} represents a rect with
\texttt{NSZeroPoint} \texttt{origin} and \texttt{NSZeroSize}
\texttt{size}.

\subsection{Geometry Functions}
The GNUstep base library provides functions and macros to do all the
most common geometry operations on \texttt{NSPoint}s, \texttt{NSSize}s
and \texttt{NSRect}s (such as determining if a point is contained in a
rectangle, computing the intersection of two rectangles, etc).  It
would be off topic to discuss them here; please refer to the base
library documentation whenever you need to do some of these common
geometrical operations.

\section{Adding a Window to your Application}
\subsection{Window Attributes}
Windows are represented in GNUstep by \texttt{NSWindow} objects.  We
are mainly interested in two attributes of a \texttt{NSWindow} object:
\begin{enumerate}
\item its \emph{content rect}, which is a \texttt{NSRect} describing
  the rectangle covered by the internal area of the window; i.e.,
  without all the decorations added by the window manager (title bar,
  border, etc);
\item its \emph{style mask}, which is an \texttt{unsigned int}
  describing the decorations of the window.  A zero style mask is also
  known as a \texttt{NSBorderlessWindowMask}, and means that the
  window has no decorations at all.  If the window needs to have
  decorations, you need to set its style mask to a combination
  (\texttt{OR}) of one or more of the following constants:
\begin{itemize}
\item \texttt{NSTitledWindowMask} = the window has a title bar;
\item \texttt{NSClosableWindowMask} = the window can be closed by the
  user (it has a close button);
\item \texttt{NSMiniaturizableWindowMask} = the window can be
  miniaturized by the user (it has a miniaturize button);
\item \texttt{NSResizableWindowMask} = the window can be resized by
  the user.
\end{itemize}
For example, a window which has a title bar, is closable and
miniaturizable will have a style mask of:
\begin{verbatim}
unsigned int windowMask = NSTitledWindowMask 
                          | NSClosableWindowMask 
                          | NSMiniaturizableWindowMask;
\end{verbatim}
where \texttt{|} is the C operator for the logical OR.
\end{enumerate}

While you might freely change the content rect of your window after
the window has been created, you can only set the style mask when the
window is first created.  You can not change the style mask of the
window after creation.

\subsection{Creating a Window}
We are now ready to show an example of creating a window: 
\begin{verbatim}
NSWindow *myWindow;
NSRect rect = NSMakeRect (100, 100, 200, 200);
unsigned int styleMask = NSTitledWindowMask 
                         | NSMiniaturizableWindowMask;

myWindow = [NSWindow alloc];
myWindow = [myWindow initWithContentRect: rect
                     styleMask: styleMask
                     backing: NSBackingStoreBuffered
                     defer: NO];
\end{verbatim}
(Please ignore the \texttt{backing:} and \texttt{defer:} arguments for
now.  The values given in the example are appropriate in most cases).
This window has an initial content rect which has its origin (its
lower left corner) positioned in the point of coordinates (100, 100)
in the screen coordinates.  In GNUstep, all coordinates system are by
default Cartesian and the origin is always in the lower left corner
(this is only the default, because inside a GNUstep window you can
then flip, rotate and (more generally) transform coordinates at your
wish).  The window has a title bar, and can be miniaturized by the
user, but not resized or closed.

\subsection{Setting the title of a window}
To set the title of the window \texttt{myWindow}, you simply do
something like:
\begin{verbatim}
[myWindow setTitle: @"This is a test window"];
\end{verbatim}

\subsection{Ordering Front a window}
Creating a window does not make it visible.  The window is ready to be
used, but it is not visible till you invoke the \texttt{orderFront:}
method, as in:
\begin{verbatim}
[myWindow orderFront: nil];
\end{verbatim}
Note that \texttt{orderFront:} takes an object argument so that it can
be invoked as an \emph{action} from a button or a menu item, but the
argument is usually ignored, which is why we simply use \texttt{nil}
for it.

If you want your window to be ordered front and to get the keyboard
focus, you do:
\begin{verbatim}
[myWindow makeKeyAndOrderFront: nil];
\end{verbatim}

\subsection{Integrating the window with your application}
We now want to show a very simple example of an application creating a
single window.  Our application will be useless if its only window is
not visible; for this reason, we do not add a close button to the
window (following the \emph{NEXTSTEP} tradition, the user needs to
select \emph{Quit} from the main menu to quit the application.  To get
the main menu, the user typically clicks the right button of his mouse
on one of the windows).

We can create the window in the same method 
\begin{verbatim}
- (void) applicationWillFinishLaunching: (NSNotification *)not;
\end{verbatim}
where we created the menu.  This method contains initialization 
code to be run \emph{before} the application is launched.

But, we can't order front (ie, make visible) windows before
\texttt{NSApp} has become active (which happens when the application
is `launched'), so that ordering front our window in
\texttt{applicationWillFinishLaunching:} would not work, because this
method is called \emph{before} the application finished launching.

To get around this problem, we implement in our application delegate
also another method, called:
\begin{verbatim}
- (void) applicationDidFinishLaunching: (NSNotification *)not;
\end{verbatim}
This method (if implemented by the application's delegate) is invoked
by the \texttt{NSApp} just after it finished launching.  We order
front our window in this method.

\subsubsection{The organization of start-up code}
The way we have organized the start-up of our application is as
follows:
\begin{verbatim}
- (void) applicationWillFinishLaunching: (NSNotification *)not;
\end{verbatim}
contains code for the initialization process of the application.  This
usually means creating the application's main objects; in our case,
the menu and the only application's window.  You can't (and shouldn't)
display any window in this method.

Instead, the method
\begin{verbatim}
- (void) applicationDidFinishLaunching: (NSNotification *)not;
\end{verbatim}
contains a sequence of user-visible actions to be taken just after the
application became active: typically, popping up a window or a panel
to the user.

It should be mentioned that this is not the only possible way of
organizing the start-up of an application.

\subsection{Source Code}
Here is the source code of our new application: 
\begin{verbatim}
#include <Foundation/Foundation.h>
#include <AppKit/AppKit.h>

@interface MyDelegate : NSObject
{
  NSWindow *myWindow;
}
- (void) createMenu;
- (void) createWindow;
- (void) applicationWillFinishLaunching: (NSNotification *)not;
- (void) applicationDidFinishLaunching: (NSNotification *)not;
@end

@implementation MyDelegate : NSObject 
- (void) dealloc
{
  RELEASE (myWindow);
}

- (void) createMenu
{
  NSMenu *menu;

  menu = AUTORELEASE ([NSMenu new]);

  [menu addItemWithTitle: @"Quit"  
        action: @selector (terminate:)  
        keyEquivalent: @"q"];

  [NSApp setMainMenu: menu];
}

- (void) createWindow
{
  NSRect rect = NSMakeRect (100, 100, 200, 200);
  unsigned int styleMask = NSTitledWindowMask 
                           | NSMiniaturizableWindowMask;


  myWindow = [NSWindow alloc];
  myWindow = [myWindow initWithContentRect: rect
                       styleMask: styleMask
                       backing: NSBackingStoreBuffered
                       defer: NO];
  [myWindow setTitle: @"This is a test window"];
}

- (void) applicationWillFinishLaunching: (NSNotification *)not
{
  [self createMenu];
  [self createWindow];
}

- (void) applicationDidFinishLaunching: (NSNotification *)not;
{
  [myWindow makeKeyAndOrderFront: nil];
}
@end

int main (int argc, const char **argv)
{ 
  [NSApplication sharedApplication];
  [NSApp setDelegate: [MyDelegate new]];

  return NSApplicationMain (argc, argv);
}
\end{verbatim}
To make the code easier to read and to understand, we have moved all
the code creating the menu into a \texttt{createMenu} method, and the
code creating the window into a \texttt{createWindow} method.

We have also implemented the \texttt{dealloc} method, which is quite
useless in this case, because we create only one \texttt{MyDelegate}
object, which is only released when the application quits (there is
few interest in releasing memory when the application is quitting,
since all memory is going to be released anyway).  But it is good
programming practice to always implement \texttt{dealloc}; this method
should release all the resources that the object was using.  In this
case, we only need to release the window object.  Please refer to the
Memory Management tutorial [yet to be written *grin*] for more
information on the \texttt{dealloc} method.

The \texttt{GNUmakefile} is the usual one:
\begin{verbatim}
include $(GNUSTEP_MAKEFILES)/common.make

APP_NAME = WindowApp
WindowApp_OBJC_FILES = MyApp.m 

include $(GNUSTEP_MAKEFILES)/application.make
\end{verbatim}

\section{Adding a Button in the Window}
We are now going to add a button inside our window.  But first, we
need to prepare the ground with a very short introduction to the
wonderful \texttt{NSView} class.

\subsection{A Quick Introduction to Views}
\texttt{NSView} is a class whose objects represent a rectangular area
in a window.  Each \texttt{NSView} has its own internal coordinate
system; \texttt{NSView} provides a complete framework for managing the
rectangle represented by the view, its coordinate system, and for
drawing and getting user events inside the view's rectangle.

Any object which can be displayed (and/or can accept user events)
inside a window (such as a button, a text field, a scroll view, a
slider, etc) is actually an instance of a subclass of \texttt{NSView}.
The specific subclass (for example, \texttt{NSButton}) implements a
specific way of drawing and managing user events inside the view's
rectangle.  An unsubclassed \texttt{NSView} by default draws nothing
and reacts to no user events, even if it contains all the machinery
(ready for subclasses to use) to do both.

When you create a window, the library automatically creates an
\texttt{NSView} covering the whole content area of the window.  This
view is called the window's \emph{content view}.  You can replace the
default content view (which is transparent and reacts to no user
events) with an \texttt{NSView} of your choice (which can be an
instance of a subclass of \texttt{NSView}, such as a \texttt{NSButton}
or a \texttt{NSScrollView}), but you should not change the rectangle
it represents, which must always cover the whole window.

If you want to have a view which covers only a part of a window, you
need to add it as a \emph{subview} of the content view.  We will
explore subviews in a forth-coming tutorial; in our first example, we
will go for the simplest possible solution avoiding use of subviews:
we will create a single \texttt{NSButton}, make a window to fit
exactly our button, and then replace the default window's content view
with our button.  In this way we'll get a window containing exactly
our button.  In the next tutorial, we'll learn how to enclose views
one over the other in a subview tree, so that we can put more things
in a single window.

\subsection{Creating a Button}
Creating a button is quite easy:
\begin{verbatim}
NSButton *myButton;

myButton = AUTORELEASE ([NSButton new]);
[myButton setTitle: @"Print Hello!"];
[myButton sizeToFit];
\end{verbatim}
\texttt{setTitle:} sets the title (string) to be displayed on the button; 
\texttt{sizeToFit} resizes the button so that it fits the title it is 
displaying.

Next, we need to set target and action of the button:
\begin{verbatim}
[myButton setTarget: self];
[myButton setAction: @selector (printHello:)];
\end{verbatim}
(\texttt{self} will refer to our custom object)

\subsection{Putting the button in the window}
We now need to modify our code to create the window so that it creates
a window exactly of the right size to contain our button.  First, we
need to get the size of the button:
\begin{verbatim}
NSSize buttonSize;

buttonSize = [myButton frame].size;
\end{verbatim}
\texttt{frame} is a method inherited by \texttt{NSView} which returns
the rectangle enclosing the button (or, more generally, the rectangle
enclosing/represented by an \texttt{NSView}).  The \texttt{origin} of
this rectangle is meaningless at this point; it is the \texttt{size}
(which was set by \texttt{sizeToFit}) what we want.

Then, we choose for the window a content rect with the \texttt{size}
of the button, and with an \texttt{origin} of our choice; (100, 100)
in this example:
\begin{verbatim}
NSRect rect;

rect = NSMakeRect (100, 100, 
                   buttonSize.width, 
                   buttonSize.height);
\end{verbatim}
At this point, we just create the window as before, but using this
rectangle (full code below).  The last thing to do is then replacing
the default window content view with our button:
\begin{verbatim}
[myWindow setContentView: myButton];
\end{verbatim}

\subsection{Source Code}
We are ready to show the whole source code (the \texttt{GNUmakefile}
is the usual one):
\begin{verbatim}
#include <Foundation/Foundation.h>
#include <AppKit/AppKit.h>

@interface MyDelegate : NSObject
{
  NSWindow *myWindow;
}
- (void) printHello: (id)sender;
- (void) createMenu;
- (void) createWindow;
- (void) applicationWillFinishLaunching: (NSNotification *)not;
- (void) applicationDidFinishLaunching: (NSNotification *)not;
@end

@implementation MyDelegate : NSObject 
- (void) dealloc
{
  RELEASE (myWindow);
}

- (void) printHello: (id)sender
{
  printf ("Hello!\n");
}

- (void) createMenu
{
  NSMenu *menu;

  menu = AUTORELEASE ([NSMenu new]);

  [menu addItemWithTitle: @"Quit"  
        action: @selector (terminate:)  
        keyEquivalent: @"q"];

  [NSApp setMainMenu: menu];
}

- (void) createWindow
{
  NSRect rect;
  unsigned int styleMask = NSTitledWindowMask 
                           | NSMiniaturizableWindowMask;
  NSButton *myButton;
  NSSize buttonSize;

  myButton = AUTORELEASE ([NSButton new]);
  [myButton setTitle: @"Print Hello!"];
  [myButton sizeToFit];
  [myButton setTarget: self];
  [myButton setAction: @selector (printHello:)];

  buttonSize = [myButton frame].size;
  rect = NSMakeRect (100, 100, 
                     buttonSize.width, 
                     buttonSize.height);

  myWindow = [NSWindow alloc];
  myWindow = [myWindow initWithContentRect: rect
                       styleMask: styleMask
                       backing: NSBackingStoreBuffered
                       defer: NO];
  [myWindow setTitle: @"This is a test window"];
  [myWindow setContentView: myButton];
}

- (void) applicationWillFinishLaunching: (NSNotification *)not
{
  [self createMenu];
  [self createWindow];
}

- (void) applicationDidFinishLaunching: (NSNotification *)not;
{
  [myWindow makeKeyAndOrderFront: nil];
}
@end

int main (int argc, const char **argv)
{ 
  [NSApplication sharedApplication];
  [NSApp setDelegate: [MyDelegate new]];

  return NSApplicationMain (argc, argv);
}
\end{verbatim}
\end{document}
